\section{The generating kernels}\label{app:kernels}
We derive the generating kernels \eqref{eq:generating-kernels}, in the notation of
Section~\ref{sec:tails}. Their order of singularity at the boundary of
the polydisk is responsible for the tangent scale. All four variables
lie in the open unit disk, so the series and their rearrangements
converge absolutely. For the pure diagonal,
\[
 D_\infty=2\sum_{p<q}(p+1)(q+1)
     (\xi^p\eta^q-\xi^q\eta^p)(\zeta^p\omega^q-\zeta^q\omega^p).
\]
Expanding, pairing the two orders $p,q$, and cancelling the equal-index
terms gives $2(A^{-2}-\Delta{}^{-2})$.

For the mixed minimum, use
\[
 \min(p+1,q+1,r+1,s+1)
   =\sum_{h\ge0}\one_{h\le p,q,r,s}.
\]
Shifting all four indices by $h$ preserves the equation $p+q=r+s$
and contributes $\Pi^h$. The unweighted total kernel is
\[
 \sum_{g\ge0}\left(\sum_{p=0}^g \xi^p\eta^{g-p}\right)
                    \left(\sum_{r=0}^g \zeta^r\omega^{g-r}\right)
     =\frac{1-\Pi}{A\Delta{}}.
\]
For distinct $\xi,\eta$ and $\zeta,\omega$, this follows by using
$\sum_{p=0}^g \xi^p\eta^{g-p}=(\xi^{g+1}-\eta^{g+1})/(\xi-\eta)$ and summing four
geometric series. Equality at repeated variables follows by continuity.
Summing the common shifts divides it by $1-\Pi$ and gives
$1/(A\Delta{})$, proving the formula for $B_\infty$.

For the gap kernel write $q=p+h$, $s=r+h$, $h\ge1$. Its wedge factors
are $(\xi\eta)^p(\eta^h-\xi^h)$ and $(\zeta\omega)^r(\omega^h-\zeta^h)$. The starting-depth sum is
\[
 \sum_{p,r\ge0}\min(p+1,r+1)(\xi\eta)^p(\zeta\omega)^r
                       =\frac1{C(1-\Pi)}.
\]
The remaining gap sum is
\[
 \sum_{h\ge1}(\eta^h-\xi^h)(\omega^h-\zeta^h)
 =\frac{\eta\omega}{1-\eta\omega}-\frac{\eta\zeta}{1-\eta\zeta}
  -\frac{\xi\omega}{1-\xi\omega}+\frac{\xi\zeta}{1-\xi\zeta}
 =\frac{(1-\Pi)(\Delta{}-A)}{A\Delta{}}.
\]
Multiplication by the cross coefficient $-2$ gives
$K_\infty=-2(\Delta{}-A)/(A\Delta{}C)$ and hence $Q_{0,\infty}=D_\infty+K_\infty$.
Finally, realignment interchanges $A$ with $C$ and leaves $\Delta{}$ fixed.
In fact the realigned term is
\[
 \calR Q_{0,\infty}=2(C^{-2}-\Delta{}^{-2})
                         -2(\Delta{}-C)/(A\Delta{}C).
\]
Adding $B_\infty+D_\infty$ leaves
\[
 2\left(\frac2{A^2}+\frac1{C^2}-\frac2{\Delta{}^2}
        -\frac1{A\Delta{}}-\frac{\Delta{}-C}{A\Delta{}C}\right).
\]
The last two fractions combine to $-1/(AC)$, giving the last line of
\eqref{eq:generating-kernels}.
